\chapter{79}

\section{Freitag, 7. September}



Joh hielt inne, senkte den Vorschlaghammer und stellte ihn vor sich ab.
Wassertropfen rannen ihm übers Gesicht. Und hätte er sie nicht gleich
noch aufgefangen, wäre seine Brille ganz von der Nase gerutscht ins Gras
gefallen. Louis hielt Joh die letzten Pfosten gerade. Bei jedem
Hammerschlag fühlte er den Boden vibrieren. Das Spannen des Drahtes
dauerte. Schliesslich stand der Zaun.

Joh schnappte nach Luft und klatschte Louis ab.

«Geschafft, unsere Schwarznasen können kommen.»

Sie packten die Werkzeuge zusammen und stiegen zur \emph{Sunnewaid}
hoch. Louis gefiel die Zusammenarbeit mit Joh. Sie arbeiteten sich
längst in die Hand. Joh war ein konzentrierter, speditiver Typ.

Sobald sie näher beisammen standen, stellte Louis ihm Fragen. Zur
\emph{Schafscheid}, zum Ablauf der \emph{Schur}, zur Geburt der Lämmer,
allesamt zum Alltag, der sich hier unten bald ändern würde. Joh schien
ob Louis’ Interesse erfreut.

Beim Wegräumen des Materials schaute sich Joh plötzlich suchend um.
Blickte zu Louis, dann hinter sich und seufzte.

«Die Drahtrolle? Haben wir die liegenlassen?»

Louis schlug sich an die Stirn.

«Ja, natürlich, die liegt noch unten am Baum.»

Louis stand auf.

«Ich hole sie.»

«Danke, ich nehm uns unterdessen ein Bier raus, okay?»

Louis rief ein \emph{Jep} zurück. Und war weg.

\sectionbreak

Beim Abstieg fragte er sich, wie das Leben hier auf der \emph{Sunnewaid}
wohl in wenigen Tagen aussah? Mit all den vielen Tieren und Menschen.
Die Betriebsamkeit würde sich bestimmt markant verdichten.

Das steile Gefälle zwang ihn automatisch in einen leichten Laufschritt.
Er wurde immer schneller. Beim Passieren einer kleinen Baumgruppe sah er
von weitem die Drahtrolle an der mächtigen Arve angelehnt stehen. Er
schnappte nach Luft. Und mit dem Stich in sein angeschlagenes Knie
wusste er es. Das Tempo hatte die Führung übernommen. Sein rechtes Bein
knickte ein. Mit einem Aufschrei torkelte er für einen kurzen Moment
nach hinten, dann nach vorn, überschlug sich und lag schliesslich
zusammengefaltet auf der Wiese.

Schon bei der kleinsten Bewegung schossen ihm Schmerztränen in die
Augen. Er liess seinen Oberkörper rückwärts ins Gras fallen, starrte vor
sich hin und fühlte sich vom Leben endgültig gehasst. Dann war ihm, als
träten seine Gedanken gegeneinander an. Als ginge es um den Gewinn eines
Pokals. Für wirkungsvollste Angstmobilisierung in Sekundenschnelle. Und
über allem lag die Frage, ob er überhaupt aufstehen konnte.

Louis stützte sich seitlich ab, zog sachte die Beine an und hob sich
langsam über die unversehrte Seite hoch. Der kleinste Druck auf das
verletzte Bein liess ihn zusammenzucken. Der Schmerz rührte ganz
offensichtlich nicht von seinem Knie. Er kam vom Oberschenkel.
Verzweifelt und mit Wucht schlug er sich die flache Hand an die Brust.
Er brauchte Hilfe. Joh. Er musste Joh anrufen. Der Griff in die leere
Hosentasche gab ihm den Rest. Natürlich. Sein Handy lag oben auf dem
Tisch vor dem Haus.

Louis hielt die Hände als Trichter vor den Mund, holte Luft und schrie
nach Leibeskräften.

«Joh! Joh!»

Er wartete. Eigentlich müsste es möglich sein, dass ihn Joh hören
konnte. Vorausgesetzt, er wäre draussen und würde keine Musik hören.

Louis atmete noch einmal tief ein und rief erneut.

«Joh!»

Er setzte zum dritten Schrei an und im selben Moment hörte er von oben
ein kräftiges

\emph{Hallo}. Johs Stimme war unverkennbar. Louis überlegte, wie er
antworten sollte. Er richtete sich gerade auf und schmetterte die Worte
zur \emph{Sunnewaid} hoch.

«Joh! Hilfe!»
